\chapter{Methodology}
\label{chap:methodology}

%% PURPOSE: This chapter explains the HOW of the research — the pipeline from raw
%% data to validated choreographies. It must justify every design choice and connect
%% each stage explicitly to an RQ. The synthetic-data section is load-bearing: it
%% must convince the reader that the pipeline can be trusted before the real
%% Pearson data is analysed.

\section{Overview}
\label{sec:methodology-overview}
% TODO: Write a short (3-5 sentence) roadmap of the chapter. Describe the four
% pipeline stages (synthetic data generation → feature engineering → pattern
% discovery → evaluation) and state which RQ each stage is primarily designed
% to address.

\section{Synthetic Data as a Methodological Choice}
\label{sec:synthetic-data-rationale}
% TODO: Justify why synthetic data was used as the primary substrate for this
% dissertation, covering:
%   (1) What it enables — known ground-truth archetypes allow quantitative
%       validation of the clustering pipeline (we can compute ARI/NMI against
%       the true labels, which is impossible with real anonymized logs).
%   (2) Practical constraints — the real Connections Academy / Pearson log data
%       is under a data-sharing agreement; synthetic data allows full open-source
%       reproducibility and public release of the pipeline.
%   (3) Generative design fidelity — explain how the synthetic generator was
%       calibrated to reproduce realistic behavioral distributions (action
%       frequencies, time-of-day patterns, session lengths) observed in the
%       K-12 online learning literature.
%   (4) Explicit scope limitations — results cannot be directly generalized to
%       real students without repeating the analysis on the real dataset; the
%       synthetic study establishes pipeline validity, not empirical claims about
%       Connections Academy students.

\section{Synthetic Data Generation}
\label{sec:synthetic-data-generation}
% TODO: Describe the synthetic data generator in detail:
%   - The five behavioral archetypes and their defining parameters:
%       Steady_Worker, Balanced_Engager, Low_Engagement, Night_Crammer,
%       Disengaged_Ghost.
%   - How students are sampled from each archetype (proportions, at-risk
%     correlation per archetype).
%   - How school-level demographic structure is assigned (32 schools, 32 states,
%     13 grade levels from K to Grade 12).
%   - How activity logs are generated per student per week (action type
%     distribution, time-of-day distribution, session length, duration).
%   - Fixed random seed (42) for full reproducibility.
%   Reference the generation script in the repository.

\section{Feature Engineering}
\label{sec:feature-engineering}
% TODO: Describe the weekly feature matrix construction (build_weekly_features.py):
%   - The unit of analysis: student-week (one row per student per calendar week).
%   - The two-block design rationale:
%       WHAT block (act__ prefix) — proportions of each semantic action type.
%       WHEN block (time__ prefix) — proportions of activity in each
%       time-of-day bin (Morning, Afternoon, Evening, Night).
%   - Why Login and Logout are excluded from the WHAT block (universal/structural
%     actions that do not discriminate behavior).
%   - Why the two blocks are kept separate rather than crossed (a crossed
%     17-action × 4-bin design would let the time dimension dominate k-means).
%   - The groupby + unstack reshaping pipeline: how the event-level log is
%     transformed into the wide student-week matrix.
%   - Metadata columns (n_sessions, active_days, total_actions) and why they
%     are excluded from the clustering feature space.

\section{Pattern Discovery (Clustering)}
\label{sec:pattern-discovery}
% TODO: Describe the clustering pipeline (run_clustering.py) in three sub-stages:
%
%   Preprocessing:
%   - L1 shape normalization (normalize_shape): converts raw action counts to
%     proportions so clustering focuses on behavioral shape, not volume.
%     Motivate with the example of two students with the same mix but different
%     total actions.
%   - Safe standardization (safe_standardize): zero-centres and scales each
%     feature to unit variance, but zeroes out near-constant columns to avoid
%     divide-by-zero issues that would corrupt k-means distances.
%
%   k-Selection:
%   - Elbow method (inertia) and Silhouette score across k = 2..10.
%   - Why silhouette is used to auto-select k (it accounts for inter-cluster
%     separation, not just intra-cluster compactness).
%   - How the two metrics can disagree and how the final k = 4 was chosen.
%
%   Clustering and Visualization:
%   - Final k-means fit with k = 4, n_init = 10, random_state = 42.
%   - UMAP projection to 2D for visual validation (n_neighbors = 15,
%     min_dist = 0.1). Clarify that UMAP is visualization-only; it does not
%     influence cluster assignments.

\section{Evaluation Design}
\label{sec:evaluation-design}
% TODO: Map each research question to its evaluation method and metric(s):
%
%   RQ1 — Identification:
%     Contingency table (cluster × archetype), purity per cluster, Adjusted
%     Rand Index (ARI), and Normalized Mutual Information (NMI) against the
%     ground-truth archetype labels from the generator.
%
%   RQ2 — Robustness:
%     (a) Temporal: week-over-week ARI and same-cluster retention %; cluster
%         transition matrix and per-cluster stickiness (diagonal).
%     (b) Cross-context: choreography share (%) distribution by grade level
%         and by state; standard deviation across contexts; cosine similarity
%         of per-cluster behavioral signatures across grade bands (K-5 vs 6-12).
%
%   RQ3 — Association with outcomes:
%     Dominant choreography per student (mode of weekly labels); one-way ANOVA
%     on Final_Grade across choreographies; eta-squared effect size; mean grade
%     stratified by at-risk status within each choreography.
%
%   RQ4 — Predictive utility:
%     Random Forest classifier trained on first-N-weeks features only (N = 4
%     and N = 8). Baseline (volume only) vs. experimental (volume + one-hot
%     dominant cluster). Primary metric: ROC-AUC. Data leakage is prevented
%     by strict temporal feature cutoff.
%
%   RQ5 — Actionability:
%     Gini importance and permutation importance from the trained Random Forest;
%     direction of effect (pass rate in top vs. bottom third of each feature);
%     SHAP (TreeExplainer) mean absolute values for corroboration.

\section{Reproducibility and Implementation Details}
\label{sec:reproducibility}
% TODO: State all implementation decisions that affect reproducibility:
%   - Fixed random seed 42 used throughout (k-means, UMAP, train/test split,
%     Random Forest, SHAP sampling).
%   - Python toolchain and versions: scikit-learn, umap-learn, pandas, numpy,
%     matplotlib, seaborn, scipy, shap (optional).
%   - Pipeline execution order:
%       generate_data.py
%       → build_weekly_features.py
%       → run_clustering.py --k 4
%       → validate_clusters.py
%       → temporal_stability.py
%       → context_stability.py
%       → outcome_association.py
%       → early_prediction.py
%       → interpretability.py
%   - All scripts located in artificial-database/ in the public repository.
%   - Hardware and runtime estimates (optional: useful for reproducibility).
